\documentclass[11pt,letterpaper]{article}
\usepackage[T1]{fontenc}
\usepackage[utf8]{inputenc}
\usepackage{lmodern,microtype}
\usepackage[margin=1in]{geometry}
\usepackage{amsmath,amssymb,amsthm,mathtools}
\usepackage{booktabs,array,longtable}
\usepackage{enumitem,listings,fancyhdr}
\usepackage{xcolor,hyperref}
\hypersetup{colorlinks=true,linkcolor=blue!45!black,citecolor=blue!45!black,urlcolor=blue!45!black,
 pdftitle={Sharp Radical Bounds and Survivor-First Compression in Exact Unknot Recognition},
 pdfauthor={Research contribution prepared for the ProveIt project}}
\usepackage[nameinlink,noabbrev]{cleveref}
\setlength{\parskip}{0.3em}
\setlength{\emergencystretch}{2em}
\allowdisplaybreaks
\pagestyle{fancy}\fancyhf{}
\fancyhead[L]{\small ProveIt / Unknot Recognition}
\fancyhead[R]{\small Radical transfer and exact survivor profiles}
\fancyfoot[C]{\thepage}
\setlength{\headheight}{14pt}
\newtheorem{theorem}{Theorem}[section]
\newtheorem{lemma}[theorem]{Lemma}
\newtheorem{proposition}[theorem]{Proposition}
\newtheorem{corollary}[theorem]{Corollary}
\theoremstyle{definition}
\newtheorem{definition}[theorem]{Definition}
\newtheorem{example}[theorem]{Example}
\newtheorem{question}[theorem]{Research question}
\theoremstyle{remark}
\newtheorem{remark}[theorem]{Remark}
\newcommand{\F}{\mathbb F_2}
\newcommand{\A}{\mathcal A}
\newcommand{\C}{\mathcal C}
\newcommand{\B}{\mathcal B}
\newcommand{\Hom}{\operatorname{Hom}}
\newcommand{\End}{\operatorname{End}}
\newcommand{\rad}{\operatorname{rad}}
\newcommand{\wt}{\operatorname{wt}}
\newcommand{\rank}{\operatorname{rank}}
\newcommand{\poly}{\operatorname{poly}}
\newcommand{\Kh}{\operatorname{Kh}}
\newcommand{\rKh}{\widetilde{\operatorname{Kh}}}
\newcommand{\id}{\operatorname{id}}
\newcommand{\code}[1]{\texttt{\detokenize{#1}}}
\lstset{basicstyle=\ttfamily\small,columns=fullflexible,keepspaces=true,
 breaklines=true,frame=single,framerule=0.3pt,showstringspaces=false}
\title{\textbf{Sharp Radical Bounds and Survivor-First\\Compression in Exact Unknot Recognition}\\[0.6em]
\large Finite homological perturbation, certified minimal models,\\
and a parameterized route to quasi-polynomial complexity}
\author{Research contribution prepared for the ProveIt project}
\date{7 October 2026}
\newcommand{\BraidGain}{2.23}
\newcommand{\ProfileLow}{0.032}
\newcommand{\ProfileHigh}{0.198}

\begin{document}
\maketitle
\begin{abstract}
We develop a chain-level complement to the earlier twist and component-compression
work in ProveIt's unknot-recognition program. For the characteristic-two arc
category on $2k$ boundary points, we give an explicit proof of the sharp radical
bound $J^{2k}=0\ne J^{2k-1}$ for $k\geq1$. The proof uses a dot derivative and a
nested-matching witness. A central marked-dot splitting further identifies the
algebra as $B\otimes\F[z]/(z^2)$, with radical index $2k-1$ for $B$.

For any bounded complex in this category, the scalar residue predicts the exact
multiplicity of every matching in every degree of a minimal representative:
$\mu_{a,h}=m_{a,h}-r_{a,h}-r_{a,h-1}$. This yields both a cheap survivor oracle and
a proof that changing pivot order cannot improve these final multiplicities for
a fixed processed tangle. A finite perturbation formula constructs the remaining
differential directly on the predicted survivors. Its maximum necessary
correction order is $2k-2$, and we exhibit complexes attaining that order.
The expensive composition count is output-sensitive in the survivor count,
subject to explicit input-validation and certificate-cost qualifications.

The package contains exact standard-library Python implementations, full
homotopy-identity checks, an independently checked braid-scanning harness, and
paired benchmarks. It checks 306 braid closures against a crossing-cube oracle
and verifies 1,365 intermediate-stage contractions. On one alternating
three-braid benchmark, the local transfer scanner is about $\BraidGain$ times faster
than the included scalar-pivot baseline; other cases are slower. These are not
benchmarks against the optimized upstream recognizer. The resulting complexity
theorem is unconditional in boundary width and minimal survivor count, but no
uniform quasi-polynomial bound for arbitrary knot diagrams is established.
\end{abstract}
\noindent\textbf{Status.} Complete written proofs of the algebraic and parameterized
statements; executable exact-arithmetic checks; no Lean formalization, no
first-in-literature priority claim, and no claim to have run the full upstream
repository test suite.
\newpage
\begingroup\small\setlength{\parskip}{0pt}\tableofcontents\endgroup
\newpage

\section{Research objective, existing work, and precise scope}
\label{sec:scope}
The objective is an exact, practically useful unknot recognizer with an eventual
worst-case running time $n^{O(\log n)}$, where $n$ measures an explicit diagram.
There are two distinct questions. One asks how to reduce the cost of the exact
Khovanov pipeline that already exists. The other asks for a global structural
argument forcing a quasi-polynomial bound. A faster numerical kernel, a smaller
intermediate complex on examples, and a global complexity theorem are different
achievements. This report treats them separately.

The inspected ProveIt subtree documents an exact but worst-case exponential
implementation. Its scanner already includes bit-packed dotted cobordisms,
compiled gluing plans, isomorphism cancellation, fill-sensitive pivot selection,
and alternative scan orders \cite{proveit}. The earlier project contribution on
twist compression replaces long local crossing blocks by smaller complexes;
the component-compression contribution accelerates individual cobordism
compositions and studies hierarchy restart counts \cite{twist,component}.
We do not reclassify those existing improvements as results of this report.

The present work addresses a different bottleneck: repeatedly performing
cobordism-valued Schur updates before knowing how many objects will survive.
It separates a differential into a readily processed scalar part and a nilpotent
remainder. The scalar part determines the final minimal object counts exactly.
The remainder can then be transferred to those objects in a uniformly bounded
number of algebraic correction stages.

\subsection{Connection with the two referenced repositories}
The methodological connection to \code{openai/math} is its finite Artin-model
construction in the section entitled ``Finite cochain models'' of the binary
algebra manuscript \cite{mathrepo}. There, residual cohomology is lifted to a
finite model using a nilpotent perturbation. We use the same general strategy,
but none of that manuscript's arithmetic conclusions is an input to our proofs.
The homological perturbation lemma itself is classical; an appropriate primary
reference is Crainic \cite{crainic}. Our category has several nonisomorphic
matching objects, rather than a single local coefficient ring, and its exact
characteristic-two radical index supplies a concrete transfer bound.

The original arc algebra and tangle formalism are due to Khovanov and Bar-Natan
\cite{khovanov,bn05,bn07}. Minimal complexes over a nilpotent radical and the
perturbation lemma are established algebraic tools. Our contribution is an
explicit, independently proved specialization, its sharpness witnesses, its
output-sensitive implementation, and its integration with the particular
recognition pipeline. We have not established bibliographical priority for the
sharp nilpotence statement or the marked splitting.

\subsection{The two obstructions that remain}
Even when the boundary is small, the number of copies of its matching objects
can be large. This distinction is essential. Kelom\"aki--Sch\"utz prove
exponential behavior for basis-expanding scanning on certain positive and
alternating three-braid families, while also giving a different polynomial-time
algorithm for three-braid Khovanov homology \cite{ks26}. Thus bounded strand
number alone cannot justify a polynomial bound on an explicitly expanded
minimal complex. Our survivor theorem makes this multiplicity obstruction
visible before the expensive elimination phase.

The separate hierarchy route also remains relevant. Lackenby's 2026 preprint
explains in Section~9 that its algorithmic steps have not been specified
precisely enough to estimate their running times \cite{lackenby}. This report
neither supplies those missing geometric routines nor derives their complexity
bounds from Khovanov calculations. Its quasi-polynomial consequence instead
has an explicit boundary-and-survivor hypothesis.

\section{The characteristic-two arc category and its intrinsic weight}
\label{sec:arc}
Fix $2k$ labeled points in a prescribed cyclic order. Let $\B_k$ be their set of
noncrossing perfect matchings. Labels carry incidence information; changing their
names does not change the algebra. For $a,b\in\B_k$, let $c(a,b)$ be the number
of circles obtained by overlaying $a$ with the reflected matching $b$.
Each circle alternates between arcs of the two matchings, with a coincident pair
counted as a two-edge circle. In particular,
\[
 1\leq c(a,b)\leq k\quad(k>0),\qquad c(a,b)=k\Longleftrightarrow a=b.
\]

We work over $\F$. After delooping, the dotted-cobordism morphism space is
\begin{equation}\label{eq:hom}
 \Hom(a,b)=\F[x_1,\ldots,x_{c(a,b)}]/(x_1^2,\ldots,x_{c(a,b)}^2)
\end{equation}
as a vector space, with its indicated monomial basis. For unequal endpoints
this notation specifies the basis, not a multiplication internal to that Hom
space. Composition is the usual saddle-gluing operation. The Frobenius algebra
on a circle is
\begin{equation}\label{eq:frob}
 A=\F[x]/(x^2),\qquad
 \Delta(1)=1\otimes x+x\otimes1,\qquad \Delta(x)=x\otimes x.
\end{equation}
A sphere without a dot is zero, a sphere with one dot is one, and positive-genus
connected components evaluate to zero in this characteristic. These are precisely
the numerical conventions used by the inspected composition routines.

The additive category of finite direct sums of these matchings is denoted
$\C_k$. Equivalently, use the finite-dimensional algebra
\[
 H_k=\bigoplus_{a,b\in\B_k}\Hom(a,b),
\]
with products zero when endpoints do not match. Our convention for $gf$ is
``$g$ after $f$.'' Quantum shifts are forgotten, while homological degrees are
retained. This is enough to compute total Khovanov ranks. The intrinsic weight
below is an algebra grading; it is not the homological degree and is not a
claim that all entries of an unshifted differential have the same weight.

\begin{definition}[Intrinsic weight]
For a basis monomial $x_S:a\to b$, define
\begin{equation}\label{eq:weight}
 \wt(x_S)=k-c(a,b)+2|S|.
\end{equation}
Let $J$ be the span of all positive-weight basis monomials in $H_k$.
\end{definition}

\begin{lemma}[Additivity of weight]\label{lem:weight}
Every nonzero monomial in the product of homogeneous morphisms of weights
$u$ and $v$ has weight $u+v$. All weights lie between $0$ and $2k$.
The weight-zero space consists exactly of the scalar identities of the
individual matchings.
\end{lemma}
\begin{proof}
Consider basis inputs in $\Hom(a,b)$ and $\Hom(b,c)$, and put
$p=c(a,b)$, $q=c(b,c)$, and $r=c(a,c)$. Glue the $p+q$ input discs along the
$k$ middle arcs. If a component has positive genus, its contribution is zero,
so suppose all surviving components have genus zero. Let $t$ be their number.
The Euler characteristic gives
\[
 p+q-k=2t-r.
\]
A genus-zero component with $s$ output boundaries contributes $s-1$ output
dots when it has no input dot, and $s$ output dots when it has one input dot.
Two or more dots on that component give zero. Consequently a surviving output
has $r-t+d$ dots, where $d$ is the total input dot count. Its weight is
\[
 k-r+2(r-t+d)=2k-p-q+2d,
\]
which is the sum of the input weights. In these arc compositions every component
has output boundary; the same Euler accounting also accommodates any scalar
sphere evaluation in a larger dotted-cobordism presentation.

Nonnegativity follows from $c(a,b)\leq k$. The largest possible weight in a Hom
space is $k+c(a,b)\leq2k$. Weight zero forces $c(a,b)=k$ and no dots, hence
$a=b$ and the identity basis element.
\end{proof}

\begin{corollary}[Residue algebra]\label{cor:residue}
The space $J$ is a nilpotent two-sided ideal, initially with the bound
$J^{2k+1}=0$, and
\begin{equation}\label{eq:quotient}
 H_k/J\cong\prod_{a\in\B_k}\F.
\end{equation}
It is the Jacobson radical of $H_k$.
\end{corollary}
\begin{proof}
The ideal and nilpotence assertions follow from weight additivity. The quotient
is the weight-zero algebra. A nilpotent ideal is contained in the Jacobson
radical, since $1-u$ has a finite geometric inverse for every element of that
ideal, even after multiplication by an arbitrary algebra element. Conversely,
a semisimple quotient has zero radical. Applying these two observations proves
the final assertion.
\end{proof}

The projection in \cref{eq:quotient} is especially inexpensive: discard every
entry between distinct matchings, and retain only the constant coefficient of
an entry between equal matchings. In the repository's packed encoding this
coefficient is the low bit. The nontrivial fact is that discarding the other
entries respects composition; \cref{lem:weight} proves exactly this fact.

\section{Dot differentiation and the sharp radical index}
\label{sec:sharp}
\subsection{A derivation in characteristic two}
For a Hom-space monomial, define
\begin{equation}\label{eq:derivative}
 \partial(x_S)=\sum_{i\in S}x_{S\setminus\{i\}}.
\end{equation}
This definition is independent of the numbering of the overlay circles.
It lowers weight by two and satisfies $\partial^2=0$.

\begin{lemma}[Leibniz rule for arc composition]\label{lem:derivation}
For composable morphisms $f,g$,
\begin{equation}\label{eq:leibniz}
 \partial(gf)=(\partial g)f+g(\partial f).
\end{equation}
\end{lemma}
\begin{proof}
On $A$, write $D(1)=0$ and $D(x)=1$. Multiplication obeys
$D\mu=\mu(D\otimes1+1\otimes D)$. The only possibly misleading case is
$x\otimes x$: the right side is $x+x=0$, as required. Similarly,
\[
 (D\otimes1+1\otimes D)\Delta(1)=0=\Delta D(1),
\]
whereas its value on $\Delta(x)$ is $1\otimes x+x\otimes1=\Delta D(x)$.
Thus both elementary saddles commute with the total dot derivative on all their
circle factors. Arc composition is a succession of these saddles, so the
identity persists under composition. Before the saddles, the total derivative
on the tensor product of the two input spaces is precisely the right side of
\cref{eq:leibniz}. No assertion of derivative naturality for arbitrary cups or
caps is needed here.
\end{proof}

\begin{theorem}[Sharp characteristic-two radical bound]\label{thm:nilpotence}
For every $k\geq1$,
\begin{equation}\label{eq:sharpnil}
 J^{2k}=0\qquad\text{and}\qquad J^{2k-1}\ne0.
\end{equation}
The nilpotence index is therefore exactly $2k$. For $k=0$, $H_0=\F$ and $J=0$.
\end{theorem}
\begin{proof}[Upper bound]
It is enough to consider products of homogeneous radical basis elements.
A product of $2k$ such factors can survive the weight bound only if every factor
has weight one and the output has weight $2k$. A weight-one monomial is
undotted, so its derivative is zero. By \cref{lem:derivation}, the product has
zero derivative.

Weight $2k$ occurs only in a diagonal Hom space, where it is spanned by the
fully dotted monomial $x_1\cdots x_k$. Its derivative is
\[
 \sum_{i=1}^k\prod_{j\ne i}x_j\ne0.
\]
Distinct summands are distinct basis monomials, including the single constant
summand when $k=1$. Thus a nonzero weight-$2k$ output cannot have zero
derivative. The proposed product vanishes. Linearity proves $J^{2k}=0$.
\end{proof}

\begin{proof}[Explicit sharpness witness]
Number the endpoints $0,1,\ldots,2k-1$. For $0\leq j\leq k-1$, set
\begin{equation}\label{eq:nested}
 a_j=\{(0,2j+1)\}\cup\{(2s-1,2s):1\leq s\leq j\}
       \cup\{(2s,2s+1):j+1\leq s\leq k-1\}.
\end{equation}
These are noncrossing matchings. Consecutive matchings differ by one saddle,
and the undotted morphism between them has weight one. Follow the path
\[
 a_0\longrightarrow a_1\longrightarrow\cdots\longrightarrow a_{k-1}
 \longrightarrow a_{k-2}\longrightarrow\cdots\longrightarrow a_0
\]
using undotted morphisms throughout. There are $2k-2$ factors.

To evaluate the round trip, first note that the forward path merges one further
circle at each step in $\Hom(a_0,a_j)$. Its value is the undotted basis element,
since $\mu(1\otimes1)=1$. Reflection gives the analogous value for the reverse
path. In their composite on the active $2j+2$ endpoints, the two input discs
are glued along $j+1$ arcs and have $j+1$ output circles. Its Euler
characteristic is $1-j$, hence its genus is zero. The resulting value is
$\Delta^{(j)}(1)$, namely the sum of products of all but one of the $j+1$
output variables. At $j=k-1$, the complete round trip is therefore
\begin{equation}\label{eq:roundtrip}
 \sum_{i=1}^k\prod_{j\ne i}x_j=\partial(x_1\cdots x_k).
\end{equation}
Multiplication by the endomorphism $x_1$ kills all but one of these summands and
produces $x_1\cdots x_k\ne0$. This is a product of $2k-1$ radical factors.
For $k=1$, the round trip has no factors and the final dot itself is the witness.
\end{proof}

\begin{remark}[Why the field matters]
The argument uses $x+x=0$ in the saddle identities. A proof over another field
must not silently reuse it. The elementary weight bound $2k+1$ is more general,
but the improvement by one here is a characteristic-two assertion.
\end{remark}

\subsection{A central marked-dot splitting}
Choose a boundary point $p$. For each matching, dot its arc incident to $p$;
the sum of these endomorphisms is denoted $z$. Sliding that dot across any
morphism shows that $z$ is central. Also $z^2=0$ and $\partial z=1$.
Set $B_k=\ker\partial\subset H_k$.

\begin{proposition}[Splitting off one square-zero variable]\label{prop:split}
For $k\geq1$, $B_k$ is a subalgebra and multiplication gives an algebra
isomorphism
\begin{equation}\label{eq:split}
 B_k\otimes_{\F}\F[z]/(z^2)\;\cong\;H_k.
\end{equation}
For any morphism $f$, its unique decomposition is
\begin{equation}\label{eq:decompose}
 f=f_0+zf_1,\qquad f_1=\partial f,\qquad f_0=f+z\partial f,
 \qquad f_0,f_1\in B_k.
\end{equation}
Moreover, if $J_B=\rad B_k$, then
\begin{equation}\label{eq:radicalsplit}
 J^r=J_B^r\oplus zJ_B^{r-1}\quad(r\geq1),\qquad
 J_B^{2k-1}=0\ne J_B^{2k-2}.
\end{equation}
For $k=1$, the last inequality reads $B_1\ne0$ and $J_B=0$.
\end{proposition}
\begin{proof}
The Leibniz rule makes $B_k$ a subalgebra. Since
$\partial(zf)=f+z\partial f$, both terms in \cref{eq:decompose} belong to
$B_k$. This identity proves existence of the splitting. If $b_0+zb_1=0$ with
$b_i\in B_k$, applying $\partial$ gives $b_1=0$ and hence $b_0=0$.
Centrality and $z^2=0$ now prove the algebra isomorphism.

The subalgebra $B_k$ is graded. It has no nonzero elements of weight $2k$ by the
fully dotted diagonal calculation. Weight $2k-1$ can only occur when
$c(a,b)=k-1$ and every overlay circle is dotted; for $k\geq2$ its derivative
is again nonzero. When $k=1$, no weight-one space exists. Thus $B_k$ has
maximum weight $2k-2$. Its positive-weight ideal is its radical by the same
argument as \cref{cor:residue}, so $J_B^{2k-1}=0$. The undotted round trip
in \cref{eq:roundtrip} belongs to $J_B^{2k-2}$ and is nonzero, proving
sharpness. Finally $J=J_B\oplus zB_k$, and an $r$-fold product contains
at most one $z$. Expanding products as subspaces gives
\cref{eq:radicalsplit}.
\end{proof}

In each individual Hom space, the splitting halves the coefficient dimension.
If the marked overlay circle is numbered zero, one may use
$y_i=x_i+x_0$ for the other circles as a basis of $\ker\partial$; these
variables also square to zero. This is a useful representation option, not a
claim that the current Python code has already realized a factor-two memory
saving. Its marked variant deliberately retains the original full encoding
for independent identity checking.

\section{The scalar residue predicts every minimal multiplicity}
\label{sec:survivors}
Let $(C,d)$ be a bounded complex in $\C_k$. Write
\[
 C^h=\bigoplus_{a\in\B_k}a^{\oplus m_{a,h}}.
\]
The scalar residue $\bar d$ is obtained by applying \cref{eq:quotient}
entrywise. It is a complex of finite-dimensional vector spaces, separately for
each matching. Let $r_{a,h}$ be the binary rank of
$\bar d^h_a:\F^{m_{a,h}}\to\F^{m_{a,h+1}}$.

\begin{definition}[Minimal complex]
A complex in $\C_k$ is minimal if every differential entry belongs to $J$.
Equivalently, its scalar residue differential is zero. This is the same stopping
condition as exhausting all invertible entries between equal matchings.
\end{definition}
The equivalence follows because
$\End(a)=\F[x_1,\ldots,x_k]/(x_i^2)$ is local: an element is a unit exactly
when its constant coefficient is one. Distinct matching objects cannot be
isomorphic, since their residues occupy different simple factors.

\begin{theorem}[Exact survivor profile]\label{thm:survivors}
Every minimal complex homotopy equivalent to $(C,d)$ has exactly
\begin{equation}\label{eq:survivors}
 \boxed{\ \mu_{a,h}=m_{a,h}-r_{a,h}-r_{a,h-1}\ }
\end{equation}
copies of matching $a$ in degree $h$. Such a minimal complex exists.
Any homotopy equivalence between two minimal complexes is a chain isomorphism.
\end{theorem}
\begin{proof}
Because the scalar projection is an algebra map, $\bar d^2=0$. Therefore the
right side of \cref{eq:survivors} is precisely the dimension of the residue
homology in that matching and degree, and is nonnegative.

Existence follows either by successively cancelling unit entries, each of which
removes two objects, or by the finite construction in \cref{sec:transfer}.
Cancellation is a chain-homotopy equivalence, and termination leaves a radical
differential. Reduce any homotopy equivalence between the original complex and
a minimal one modulo $J$. It remains a homotopy equivalence of vector-space
complexes, while the minimal target has zero differential. Its term dimensions
must equal the residue homology dimensions, proving the formula.

Now suppose $f$ is a homotopy equivalence between two minimal complexes.
Modulo $J$ all differentials vanish, so a homotopy inverse is an actual inverse
to $\bar f$ in every degree. The scalar matrix $\bar f$ is invertible within
each matching block. Lift its inverse and factor $f$ as its scalar invertible
part times $1+N$, where $N$ is a matrix over $J$. The matrix ideal has the
same nilpotence bound as $J$, since each entry of a long enough matrix product
is a sum of vanishing products of radical elements. A finite geometric series
inverts $1+N$. Thus $f$ is invertible in every degree; its inverse is a chain
map because $f$ was one.
\end{proof}

\subsection{Consequences for the existing cancellation strategy}
For a fixed processed tangle, exact minimal multiplicities do not depend on the
order of scalar-unit cancellations. Min-fill, LIFO, or another complete pivot
strategy can change intermediate fill, coefficient density, and runtime, but
not the final counts in \cref{eq:survivors}. This is a stronger statement than
agreement of final closed homology ranks: it applies before closure, matching
by matching and degree by degree.

The formula is also a resource oracle. A caller can compute
$R=\sum_{a,h}\mu_{a,h}$ before cobordism-valued Schur updates and decline an
explicit representation whose minimal survivor count already exceeds a budget.
That refusal is \emph{not a knot verdict}. In particular, a large partial-tangle
profile does not imply that its eventual closure is nontrivial.

The lower bound is representation-specific. It applies to finite direct sums of
these matching objects up to chain homotopy. It does not rule out circuits,
recurrences, macro complexes, compressed multiplicities, derived descriptions
using different generators, or a decision procedure that avoids computing all
Khovanov homology. This qualification is precisely where future
quasi-polynomial progress may enter.

\section{Finite transfer directly to the survivors}
\label{sec:transfer}
\subsection{Constructing the scalar contraction}
Let $d_0$ be the scalar lift of $\bar d$: equal-matching constant coefficients
are multiplied by the corresponding identity cobordisms, and other entries
are zero. These scalar blocks form a split subalgebra, so $d_0^2=0$. Put
\[
 \delta=d-d_0=d+d_0.
\]
Every entry of $\delta$ belongs to $J$.

Binary linear algebra decomposes each matching block of $(C,d_0)$ as homology
plus contractible two-term pairs. Let $V$ denote the homology objects with the
multiplicities \cref{eq:survivors}. Choose maps
\[
 i:V\to C,\qquad p:C\to V,\qquad h:C^h\to C^{h-1}
\]
whose entries are scalar identities and which satisfy
\begin{equation}\label{eq:scalarcontraction}
 pi=1,\quad d_0i=0,\quad pd_0=0,\quad
 1+ip=d_0h+hd_0,\quad h^2=hi=ph=0.
\end{equation}
One explicit construction chooses, in every degree, a basis $B$ of boundaries,
a complement $H$ to $B$ in the cycles, and lifts $L$ of the next boundaries.
In the basis $(B,H,L)$, let $i$ include $H$, let $p$ project to it, and let $h$
map a boundary to its selected lift in the preceding degree. The implementation
builds these bases with tracked binary column elimination.

\begin{theorem}[Sharp finite transfer]\label{thm:transfer}
For $k\geq1$, set
\begin{align}
 I&=\sum_{r=0}^{2k-1}(h\delta)^r i,&
 P&=p\sum_{r=0}^{2k-1}(\delta h)^r,\label{eq:IP}\\
 H&=\sum_{r=0}^{2k-1}(h\delta)^r h,&
 D&=\sum_{r=0}^{2k-2}p\delta(h\delta)^r i.\label{eq:HD}
\end{align}
Then $(V,D)$ is minimal and the following identities hold:
\begin{gather}\label{eq:certificate}
 D^2=0,\quad dI=ID,\quad Pd=DP,\quad PI=1,\quad
 1+IP=dH+Hd,\\
 H^2=0,\qquad HI=0,\qquad PH=0.\nonumber
\end{gather}
Consequently it is homotopy equivalent to $(C,d)$. For $k=0$, simply take
$D=0$ and the scalar contraction.
\end{theorem}
\begin{proof}
Matrices with entries in $J$ form a nilpotent ideal of index at most $2k$,
independently of matrix size. Since $h$ is scalar, $h\delta$ and $\delta h$
are in this ideal. Their geometric series therefore terminate as displayed.
Each summand of $D$ has at least one radical factor, proving minimality.
The omitted possible last summand has $2k$ radical factors and vanishes.

For completeness, the algebra behind the perturbation step can be recorded
without an infinite-series argument. Define
\[
 Q=\delta(1+h\delta)^{-1}=(1+\delta h)^{-1}\delta.
\]
Finite multiplication gives
\[
 Q=\delta+\delta hQ=\delta+Qh\delta,
 \qquad d_0Q+Qd_0+QipQ=0.
\]
The last identity follows by substituting $ip=1+d_0h+hd_0$ and
$d_0\delta+\delta d_0+\delta^2=0$, then multiplying by the two finite
inverses. In these terms,
\[
 D=pQi,\quad I=i+hQi,\quad P=p+pQh,\quad H=h+hQh.
\]
For example,
$D^2=p(QipQ)i=p(d_0Q+Qd_0)i=0$.
Expanding $dI+ID$ and using the same three identities cancels all terms;
the dual calculation gives $Pd+DP=0$. Expansion of $dH+Hd$ gives
$1+IP$. The side conditions $h^2=hi=ph=0$ immediately imply
$PI=1$, $H^2=0$, $HI=0$, and $PH=0$ from the displayed formulas.
These are the finite special-deformation-retract formulas of the classical
perturbation lemma \cite{crainic}, here justified by the explicit radical bound.
\end{proof}

\subsection{Why the stated correction depth cannot be reduced}
\begin{proposition}[Sharp transfer order]\label{prop:sharptransfer}
For each $k\geq1$, there is a two-term complex with $4k-2$ objects and exactly
two residue survivors for which the first nonzero transferred-differential
summand is
\[
 p\delta(h\delta)^{2k-2}i.
\]
Thus $2k-2$ is the best possible universal correction-order bound for this
category.
\end{proposition}
\begin{proof}
Take the $q=2k-1$ radical factors $f_1,\ldots,f_q$ in the sharpness witness
of \cref{thm:nilpotence}, on a path of matching objects
$a^{(0)},\ldots,a^{(q)}$. Place a survivor $U=a^{(0)}$ in degree zero and a
survivor $V=a^{(q)}$ in degree one. For each internal path vertex add a
contractible pair $B_j\xrightarrow{1}C_j$, with $B_j$ in degree zero and
$C_j$ in degree one, both of matching type $a^{(j)}$.

Let $d_0$ consist of these identity arrows. Define $\delta$ by
\[
 U\xrightarrow{f_1}C_1,\quad
 B_1\xrightarrow{f_2}C_2,\quad\ldots,\quad
 B_{q-1}\xrightarrow{f_q}V.
\]
All maps go from degree zero to degree one, so $d^2=0$ automatically.
The scalar contraction sends $C_j$ back to $B_j$. The only route from $U$
to $V$ in the transfer series uses all $q$ radical factors and $q-1$ copies of
$h$. Its value is the nonzero fully dotted endomorphism. Every shorter route
ends in a contractible target and is killed by $p$. The object count is
$2+2(q-1)=4k-2$.
\end{proof}

This witness is an algebraic complex in the tangle category. We do not assert
that every such witness is the scanner complex of a knot diagram. The theorem
establishes sharpness for the algebraic transfer routine, not a topological
worst-case lower bound for all recognizers.

\subsection{A block-matrix pitfall}
A scalar unit of $\End(a)$ is an involution over $\F$, since every element
with zero constant term squares to zero. This existing scalar shortcut must
not be used for matrices. In $\F[x,y]/(x^2,y^2)$, let
\[
 N=\begin{pmatrix}0&x\\y&0\end{pmatrix}.
\]
Then $N^2=xyI\ne0$ and $N^3=0$, so
\[
 (I+N)^{-1}=I+N+N^2\ne I+N.
\]
This example is included in the executable tests. When a full block inverse is
needed, geometric doubling is valid: powers of the \emph{same} matrix commute,
and
\[
 \prod_{j=0}^{s-1}(I+N^{2^j})=\sum_{r=0}^{2^s-1}N^r.
\]
Choose $2^s$ at least the nilpotence bound. No entrywise Frobenius rule for
noncommuting matrix terms is being assumed.

\subsection{Two-pass transfer using the marked splitting}
\begin{proposition}[One final square-zero correction]\label{prop:twopass}
Write a differential uniquely as
\[
 d=d_B+ze,\qquad d_B=d+z\partial d,\qquad e=\partial d,
\]
with entries of $d_B,e$ in $B_k$. Then
$d_B^2=0$ and $d_Be+ed_B=0$. Suppose
$(I_B,P_B,H_B)$ contracts $(C,d_B)$ to its minimal complex $(V,D_B)$.
A minimal differential for $d$ is
\begin{equation}\label{eq:twopass}
 D=D_B+zP_BeI_B.
\end{equation}
The corresponding contraction is
\begin{align*}
 I&=I_B+zH_BeI_B,& P&=P_B+zP_BeH_B,\\
 H&=H_B+zH_BeH_B.
\end{align*}
\end{proposition}
\begin{proof}
The direct sum decomposition $H_k=B_k\oplus zB_k$ and $z^2=0$ separate the
two coefficients of $d^2$: they are $d_B^2$ and $d_Be+ed_B$.
Apply the perturbation lemma a second time, now with perturbation $ze$.
Every product containing two such perturbations vanishes because $z$ is
central and $z^2=0$. Thus only the displayed correction remains.
The scalar residues of $d$ and $d_B$ coincide, so their survivor objects
coincide as well. Finally, $D_B$ is radical and the extra term is divisible by
$z$, hence the final differential is minimal.
\end{proof}

The first pass may use the smaller radical index $2k-1$ in $B_k$. A practical
reduced-basis implementation would need only half the coefficients in each
nonzero Hom space. Our supplied two-pass variant is instead a proof-oriented
full-encoding implementation; its 308 complete contraction checks are reported
below, but it is not included in the speedup claims.


\section{Output-sensitive cost and the quasi-polynomial parameter regime}
\label{sec:cost}
\subsection{Computational model and input qualifications}
Let $M$ be the number of objects in a supplied complex before reduction, and
let $R=\sum_{a,h}\mu_{a,h}$ be the exact survivor count. A morphism in a
$2k$-point category has at most $2^k$ coefficients. Let $\A(k)$ bound the
bit cost of composing or adding two such morphisms, including the required
topological-plan work. Standard finite basis arithmetic gives
\[
 \A(k)=\poly(k)2^{O(k)}.
\]
This loose bound is enough for the asymptotic result. A compatible faster
component-compression kernel can improve the exponential constant, but is not
required for the proofs or the experiments here.

All complexity statements charge for packed bit strings, not merely for
unit-cost arithmetic on arbitrarily large Python integers. Index and endpoint
bookkeeping contributes logarithmic factors in the explicit input size and
$M$; we suppress these by $\widetilde O$. A complex is assumed either to be
known valid from its construction or to have been separately checked.
Verifying $d^2=0$ by generic dense multiplication can itself cost
$\widetilde O(M^3\A(k))$. That cost must be added when requested.

\begin{theorem}[Survivor-first construction cost]\label{thm:cost}
On a valid input differential, a minimal differential and its inclusion map
can be constructed in
\begin{equation}\label{eq:cost}
 \widetilde O\!\left(M^3+M^2(1+kR)\A(k)\right)
\end{equation}
bit operations. The $M^3$ term is scalar binary linear algebra. The expensive
cobordism-composition count in the transfer phase is $O((k+1)M^2R)$.
A complete explicit contraction, including its $M\times M$ homotopy, can
instead be produced and checked within
\begin{equation}\label{eq:certificatecost}
 \widetilde O\!\left(M^3+(k+1)M^3\A(k)\right).
\end{equation}
\end{theorem}
\begin{proof}
Build the scalar residue, compute its ranks, and construct the $(B,H,L)$ bases
from \cref{sec:transfer}. Dense binary elimination and basis inversion cost
$O(M^3)$ bit operations conservatively, after summing over the matching-degree
blocks. Reading and forming the radical coefficient arrays is covered by
$\widetilde O(M^2\A(k))$.

To construct $I$, start with the $M\times R$ scalar matrix $i$ and iterate
\begin{equation}\label{eq:iteration}
 T_0=i,\qquad T_{r+1}=h(\delta T_r),\qquad I=\sum_rT_r.
\end{equation}
Do not form the $M\times M$ matrix $h\delta$ first. Each application of
$\delta$ to $T_r$ uses at most $M^2R$ morphism compositions. Applying the
scalar matrix $h$ only sums coefficients, with no new surface composition.
There are at most $2k-1$ noninitial stages. Finally compute $D=p\delta I$.
This proves \cref{eq:cost}; its extra input term also covers the case $R=0$.
When $k=0$, the radical remainder is zero, so the construction is purely scalar.

For $P$, the analogous row iteration uses at most $O(kM^2R)$ compositions.
Constructing $H$ explicitly involves $M$ columns rather than $R$, giving the
$O(kM^3)$ bound. Each of the finitely many certificate identities can be
checked by generic matrix multiplication in $O(M^3\A(k))$ time. This proves
\cref{eq:certificatecost}.
\end{proof}

There are three distinct modes. The profile oracle performs only the binary
rank calculation. The production transfer path can skip generic input-square
verification when validity follows from the crossing construction. The
proof-oriented path computes a full contraction and checks it. The code's
\code{transfer(..., validate=False)} implements the second mode; its safer
standalone default validates the input. The output-sensitive bound must not be
advertised for that default without adding its validation cost.

Working coefficient matrices require $O(M^2 2^k)$ bits for the input and
$O(MR2^k+R^2 2^k)$ bits for transfer columns and output, in addition to scalar
basis data and index storage. This is an accounting of matrix payloads, not a
tight bound on the shipped Python process: its geometry caches and object
headers are additional. Discarding or bounding caches preserves the time bound.
A conservative overall space bound is the stated time bound. A production
memory claim should separately account for cache policy and representation
metadata.

\subsection{A comparison with scalar-unit cancellation}
Generic dense cancellation can require $O((M-R)M^2)$ cobordism-valued Schur
updates. The transfer construction replaces this coarse count by
$O(kM^2R)$ expensive compositions for $k>0$, plus binary work. The promising
regime is therefore large cancellation, small $R$, and relatively expensive
morphism arithmetic. There is no theorem saying it wins on every sparse
instance: a good pivot sequence may make almost no fill, whereas scalar basis
changes can be dense. An adaptive implementation should retain both paths.

If the radical remainder has minimum positive intrinsic weight $s>1$, then a
product of more than $\lfloor2k/s\rfloor$ such factors vanishes already by
weight. This yields a further safe truncation of \cref{eq:iteration}; endpoint
weights can refine it more. The supplied main implementation uses the uniform
sharp bound and early termination on a zero iterate. Weight-aware truncation
is an additional implementation opportunity, not part of the reported timings.

\subsection{A complete parameterized recognition theorem}
Let $N\geq2$ be the explicit input size, including strands when a braid is
supplied. Consider a \emph{valid supplied arc-category scan}: every interface
is modeled by some $\C_{k_t}$, and the crossing order and local attachments
are specified, not found by an uncharged search. Braid scanning in
\cref{sec:braid} is one concrete instance. Define
\begin{equation}\label{eq:parameters}
 k_* = \max_t k_t,\qquad
 R_* = \max\bigl(1,\text{number of objects after minimal reduction at stage }t\bigr).
\end{equation}

\begin{lemma}[Constant local expansion]\label{lem:expansion}
Attaching one crossing to a circle-free matching creates at most two new closed
circles in either smoothing. Consequently the number of objects immediately
before the next elimination is at most eight times the preceding survivor
count. For the braid attachment used in the harness, the bound is three times.
\end{lemma}
\begin{proof}
Every new closed circle must contain at least one of the two newly attached
smoothing arcs, since the old object has no closed circles. There are at most
two such circles. Delooping $c$ circles produces $2^c\leq4$ copies in each
of the two smoothings, for a total of at most eight. In a braid attachment the
identity smoothing creates no circle, and the turnback smoothing creates at
most one, giving $1+2=3$.
\end{proof}

\begin{theorem}[Parameterized recognition bound]\label{thm:qp}
For a one-component diagram with a valid supplied scan as above, crossing
attachment, survivor-first reduction, and exact closure give an exact unknot
recognition algorithm. There is a fixed polynomial $p$ for which its running
time is bounded by
\begin{equation}\label{eq:scanbound}
 T\leq p(N,k_*)\,2^{O(k_*)}R_*^3.
\end{equation}
The constants do not depend on the input. In particular,
\begin{align}
 k_*+\log_2(R_*+1)&=O(\log N)
       &&\Longrightarrow\quad T=N^{O(1)},\label{eq:polyregime}\\
 k_*+\log_2(R_*+1)&=O((\log N)^2)
       &&\Longrightarrow\quad T=N^{O(\log N)}.\label{eq:qpregime}
\end{align}
\end{theorem}
\begin{proof}
At each stage, \cref{lem:expansion} gives $M\leq8R_*$. Local crossing maps
can be assembled with $\poly(N,k_*)2^{O(k_*)}R_*^2$ work by finite-basis
arithmetic. Apply \cref{thm:cost}, or even its full-certificate bound, at each
of at most polynomially many elementary stages. Starting with the identity or
empty matching, the same expansion bound ensures $\log R_*=\poly(N)$,
so suppressed index costs are absorbed in $p$. The resulting cost is bounded
as in \cref{eq:scanbound}.

A scan that ends with empty boundary already has a vector-space complex.
For an open braid scan, closing a matching $a$ by the fixed crossingless
closure $e$ yields $c(e,a)\leq k_*$ circles, so the explicit closed complex
has at most $2^{k_*}R_*$ basis vectors. Binary Gaussian elimination costs at
most the cube of this number, still within \cref{eq:scanbound}. Correctness
follows from local chain-homotopy equivalence and the rank criterion proved
in \cref{lem:detect} below. Taking logarithms of the bound gives both
parameter-regime consequences.
\end{proof}

\paragraph{What this theorem does not charge for, and does not prove.}
It does not search over all crossing orders or supply a topological proof that
an arbitrary knot admits a scan satisfying \cref{eq:qpregime}. The profile is
an exact diagnostic once a stage differential has been constructed; it is not
an independently available succinct certificate for the entire input. Thus
\cref{thm:qp} is an unconditional parameterized theorem and a conditional
route to a \emph{uniform} quasi-polynomial algorithm. The latter structural
step remains open in this work.

\section{An exact braid-scanning validation harness}
\label{sec:braid}
The package includes a small complete algebraic harness to exercise the new
reduction on genuine tangle complexes. It is deliberately separate from the
optimized upstream \code{FastScan} representation.

Start with the identity matching between the top and bottom endpoints of a
$b$-strand braid. For each braid generator attach its two-resolution complex
at the relevant bottom endpoints, deloop any new circle, and reduce using
either the transfer routine or an independent scalar-pivot baseline.
Positive and negative generators use opposite choices of the two smoothings;
the final homological degrees are shifted by the negative-crossing count.
The geometry of the crossing and delooping plans is adapted from the inspected
repository \code{geometry.py} and \code{algebra.py}. Static composition uses
a previously archived numerical excerpt of \code{planar.py}.

For closure, use the representable functor
\[
 \Hom(e,-):\C_b\longrightarrow\text{finite-dimensional }\F\text{-vector spaces},
\]
where $e$ pairs each final bottom endpoint with its original top endpoint.
A basis for $\Hom(e,a)$ consists of dot monomials on the closed smoothing
circles, and a differential $f:a\to a'$ acts by $g\mapsto fg$.
This is the ordinary Frobenius evaluation of that crossingless closure.
As an additive functor it preserves every chain-homotopy identity, so reducing
the open complex before closing does not change the resulting homology.

\begin{lemma}[The rank criterion used for knots]\label{lem:detect}
For a knot $K$,
\[
 K\text{ is the unknot}\quad\Longleftrightarrow\quad
 \dim_{\F}\Kh(K;\F)=2
 \quad\Longleftrightarrow\quad\dim_{\F}\rKh(K;\F)=1.
\]
\end{lemma}
\begin{proof}
Here is the characteristic-two relationship between the reduced and unreduced
complexes. On each closed smoothing take the total dot derivative
$\partial$; by the saddle calculation in \cref{lem:derivation} it commutes
with the Khovanov differential. Let $X$ be multiplication by a dot on a chosen
marked component. It is also a chain map, $X^2=0$, and
$\partial X+X\partial=1$. Hence the unreduced complex splits into two
isomorphic copies of $\ker\partial$, one multiplied by $X$, and the reduced
complex can be identified with $\operatorname{im}X=X\ker\partial$.
For homological-degree counts, the two copies therefore have equal homology.
Quantum shifts, which differ between the copies, are irrelevant here.

The unknot has reduced integral Khovanov homology $\mathbb Z$. Conversely,
mod-two reduced rank one forces rational reduced rank at most one by the
universal coefficient theorem. Rational reduced Khovanov homology of a knot
has positive rank: its graded Euler characteristic is the normalized Jones
polynomial, whose value at its normalization point is one. Its rational rank
is therefore one. The unknot-detection theorem of Kronheimer--Mrowka then
implies that $K$ is the unknot \cite{km}.
\end{proof}

For links, the harness still computes unreduced homology, but the lemma is not
used to make a knot verdict. A recognition wrapper must first check that the
input has one component. For a braid closure this is the cycle count of the
braid permutation. No resource limit or partial profile is substituted for
that topological check or for a completed homology computation.

\section{Reproducible computations and performance measurements}
\label{sec:experiments}
All reported calculations use exact integers and binary arithmetic. The Python
standard library is sufficient; no floating-point arithmetic enters a theorem
check, a rank, or a differential. Floating-point values below are wall-clock
measurements only. Timings were obtained in this container with CPython 3.13.5;
the full version and platform string, individual repetitions, and operation
counts are recorded in \code{results/benchmark.json}.

\subsection{Algebra, minimality, and homotopy certificates}
The following checks were completed in this session.
\begin{center}\small
\begin{tabular}{p{0.67\linewidth}r}
\toprule
Check & Count\\\midrule
Exhaustive basis-monomial products for $1\leq k\leq3$ & 2,244\\
Leibniz and homogeneous-weight checks on those products & 2,244 each\\
Nonzero maximal radical words, $1\leq k\leq12$ & 12\\
Sharp-depth full transfer certificates, $1\leq k\leq8$ & 8\\
Random conjugated-complex full transfer certificates & 500\\
Independent scalar-pivot minimal-profile comparisons & 500\\
Marked two-pass full contraction certificates & 308\\
Unit-test methods, including invalid input and matrix pitfall & 18\\
\bottomrule
\end{tabular}
\end{center}
The 500 random complexes start with a sum of contractible pairs and a minimal
residual block, then undergo degree-preserving invertible elementary basis
changes. Thus validity is built into their construction, not guessed from a
random matrix. The verifier nevertheless checks the square of the differential,
the square of the transferred differential, both chain-map identities, $PI=1$,
the homotopy identity, and all three side conditions. The last check confirms
that the output differential has zero scalar residue.

For additional finite evidence, the script computes the entire radical-power
filtration for $k\leq4$ by linear spans of products. Here the first entry is
$\dim H_k$, followed by $\dim J,\dim J^2,\ldots$:
\begin{center}\small
\begin{tabular}{cl}
\toprule
$k$ & Dimensions of successive radical powers, including $J^0=H_k$\\\midrule
1 & $2,\ 1,\ 0$\\
2 & $12,\ 10,\ 6,\ 2,\ 0$\\
3 & $104,\ 99,\ 82,\ 52,\ 22,\ 5,\ 0$\\
4 & $1092,\ 1078,\ 1008,\ 826,\ 546,\ 266,\ 84,\ 14,\ 0$\\
\bottomrule
\end{tabular}
\end{center}
These data agree with the sharp index, but the proof of
\cref{thm:nilpotence} is independent of finite testing. Similarly, the tested
sharp transfer depths $0,2,4,\ldots,14$ are not the basis of the general
sharpness proof.

\subsection{Independent closed-braid comparisons}
The braid test set consists of six explicit examples and 100 seeded random
words for each of two, three, and four strands. Random words have between
zero and nine crossings; the fixed seed is \code{202610072}. The total is
306 closures, including both knots and links. For each case, both the local
pivot scanner and the transfer scanner are compared degree by degree against
an independent reduced crossing-cube implementation from the prior twist
package. The expected unreduced profile is twice the reduced profile, by the
chain splitting used in \cref{lem:detect}.

All comparisons passed. The transfer scanner also produced and verified a full
contraction at every nonempty stage: 1,365 stage certificates in total. The two
scanners agreed on every intermediate survivor count. Familiar calibration
cases include the one-strand unknot with unreduced rank two, the closure of
$\sigma_1^3$ with rank six, and the closure of
$(\sigma_1\sigma_2^{-1})^2$ with rank ten.

The independence is limited and explicitly documented. The crossing-cube
oracle assembles smoothing graphs and differential matrices independently and
does not use the arc-transfer formulas. However, the static arc-composition
excerpt and the crossing-plan conventions originate in the same project as
the target scanner. We have not run the full upstream scanner, Rust port, or
their complete regression suites in this environment.

\subsection{Paired timings with favorable and unfavorable cases}
Each timing is the median of five repetitions. Method order alternates, and
geometry caches are fresh for each timed run. For fixed-complex tests, input
construction and validation are outside the measured reduction. For braid
scans, the two methods share the same local geometry harness; generic
input-square validation and full certificate construction are disabled during
timing for both. Exact output checks are performed separately. A ratio above
one means transfer is faster than the included pivot baseline.

\begin{center}\small
\begin{tabular}{lrrrrr}
\toprule
Fixed complex & $M$ & $R$ & Transfer (ms) & Pivot (ms) & Ratio\\\midrule
Gauge, $k=3$ & 34 & 2 & 2.167 & 1.951 & 0.90\\
Gauge, $k=3$ & 66 & 2 & 2.799 & 9.072 & 3.24\\
Gauge, $k=4$ & 68 & 4 & 5.077 & 28.659 & 5.65\\
Gauge, $k=3$ & 100 & 4 & 24.242 & 28.049 & 1.16\\
Sharp zigzag, $k=8$ & 30 & 2 & 0.673 & 0.481 & 0.71\\
\bottomrule
\end{tabular}
\end{center}

On these five fixed complexes the profile oracle alone took approximately
$\ProfileLow$--$\ProfileHigh$ milliseconds. This is a particularly safe
integration target: it is inexpensive relative to either full reduction path
and gives an exact answer about minimal multiplicities.

\begin{center}\small
\begin{tabular}{lrrrrr}
\toprule
Braid word & $M$ & $R$ & Transfer (ms) & Pivot (ms) & Ratio\\\midrule
$\sigma_1^{25}$ & 74 & 26 & 22.893 & 19.382 & 0.85\\
$(\sigma_1\sigma_2)^6$ & 39 & 17 & 18.751 & 14.028 & 0.75\\
$(\sigma_1\sigma_2^{-1})^6$ & 895 & 609 & 86.528 & 192.756 & 2.23\\
Braid-relator unknot & 12 & 5 & 12.905 & 9.975 & 0.77\\
\bottomrule
\end{tabular}
\end{center}

For the braid table, $M$ and $R$ are maxima across the scan. The last word is
\[
 (\sigma_1\sigma_2\sigma_1\sigma_2^{-1}\sigma_1^{-1}\sigma_2^{-1})^4
 \sigma_1\sigma_2.
\]
The bracketed word is trivial by the braid relation, and the final closure is
an unknot. Its test result is rank two. It is a useful warning against selecting
only nontrivial knots or only dense algebraic matrices as performance evidence.

These measurements support an adaptive design, not unconditional replacement.
The favorable alternating example has genuine tangle multiplicity growth; the
sparse sharpness witness and several simple braids favor scalar cancellation.
Moreover, an upstream recognition pipeline would reject many of the displayed
nontrivial knots using cheaper filters, before reaching Khovanov homology.
The measurements therefore do \emph{not} establish a full-recognizer speedup or
a speedup over the optimized upstream \code{FastScan}. Millisecond timings also
have ordinary platform noise; all individual repetitions are retained so these
ratios can be re-evaluated rather than treated as asymptotic evidence.

\section{Integration strategy and correctness boundaries}
\label{sec:integration}
The proposed destination is a new subtree such as
\path{Topology/UnknotRecognition/research/radical_transfer/}.
No current defaults should change solely on the basis of this report.

The lowest-risk first step is the non-mutating survivor oracle. The included
adapter accepts the audited bit-valued reference \code{ScanComplex} interface,
not the distinct optimized \code{FastScan} layout. It reads matching types,
homological degrees, and equal-matching low bits. It neither composes
cobordisms nor rewrites the upstream object graph. A stage instrumented before
and after existing elimination should have exactly the predicted number of
survivors. The adapter has a mock-interface test but has not been run against
a complete upstream checkout; that remains an integration check.

The second step is an opt-in transfer backend with a tested conversion between
matching IDs, circle-basis ordering, and sparse differential matrices. A full
contraction verifier should be enabled on small development cases. The third
step is paired profiling on inputs that actually reach the backend after the
existing filters, including hard unknots and negative controls. Only those
measurements can justify an adaptive production default.

\paragraph{Preserve the differential, not merely the counts.}
The profile formula does not say that the radical remainder can be discarded.
An output with the correct objects and the zero differential will generally be
wrong. The finite transfer terms are essential; \cref{prop:sharptransfer}
exhibits a surviving map that appears only at the maximum allowed order.

\paragraph{Preserve scope when interpreting a lower bound.}
Pivot-order invariance concerns the same processed tangle and the same category.
Different scan orders give different prefixes. Twist compression and more
succinct objects change the representation. Neither is contradicted by the
minimality theorem.

\paragraph{Preserve budget and input semantics.}
A negative predicted multiplicity signals an invalid residue complex or an
implementation error, not a knot property. A survivor ceiling produces a
resource outcome, never ``knotted.'' A rank criterion is applied only to a
completed calculation and a verified one-component input. The simple Python
harness has no certified production memory ceiling; its cache and allocation
behavior must not be advertised as one.

\paragraph{Keep the verification claims at their actual level.}
Written proofs establish the stated algebraic theorems. Exact computations
check finite instances and implementation identities. They are not a Lean
formalization, an independent proof of every topological convention in a large
codebase, or a certificate that the unrestricted asymptotic goal has been met.
The archive includes seeds, raw timing repetitions, source identifiers, and a
reproduction script to make further review concrete.

\section{Further research questions and specific proof targets}
\label{sec:questions}
\begin{question}[A cheaply enumerable family of favorable unknot models]
Can one construct, in quasi-polynomial time, a quasi-polynomial-sized family of
verified diagram transformations and scan orders such that every unknot input
has at least one member satisfying \cref{eq:qpregime}? The transformations
must preserve the represented knot, and their encoding and execution costs
must be included. Merely knowing that a favorable presentation exists is not
enough.
\end{question}
This target can be weaker than controlling the profile of every knot and every
scan. A precise sufficient condition is useful:
\begin{proposition}[A one-sided model-cover criterion]\label{prop:cover}
Suppose an algorithm constructs and verifies, in $N^{O(\log N)}$ time,
at most $N^{O(\log N)}$ equivalent
scan inputs, each of size at most $N^{O(\log N)}$. Suppose a fixed bound
$T_0(N)=N^{O(\log N)}$ guarantees that, whenever the original input is the
unknot, at least one of the resulting exact recognizer runs finishes within
$T_0(N)$. Then bounded simulation of these runs decides unknot recognition in
$N^{O(\log N)}$ time.
\end{proposition}
\begin{proof}
Run each recognizer for $T_0(N)$ steps. Accept if any completed run identifies
the unknot. Equivalent-input verification and exactness make acceptance sound.
If none accepts, the stated covering guarantee rules out an unknot input.
The product of the family-size bound and the per-run bound remains
$N^{O(\log N)}$. Timed-out individual runs are not assigned verdicts; the
negative conclusion uses the global covering theorem.
\end{proof}
The proposition does not provide the required cover. Its purpose is to specify
an exact topological/combinatorial target without incorrectly demanding that
all nontrivial knots have small Khovanov complexes.

\begin{question}[Succinct minimal models beyond explicit multiplicities]
Can the degree-indexed matching multiplicities and radical differentials be
stored by circuits or recurrences whose size is quasi-polynomial even when
$R_*$ is exponential? A useful representation must support crossing attachment,
composition, finite transfer, closure, and exact rank-one testing within proved
bit costs. Integer-encoding of a large multiplicity alone does not provide the
corresponding differential or a fast rank algorithm.
\end{question}
The three-braid distinction between exponentially large Betti numbers and
polynomial computation in a more structured description \cite{ks26} is a
concrete reason to pursue this direction. The exact survivor theorem identifies
what cannot be removed by ordinary pivoting, and therefore what must be
represented more succinctly.

\begin{question}[A practical reduced-basis marked backend]
Implement the algebra $B_k$ in its $2^{c(a,b)-1}$-dimensional Hom bases and
combine it with \cref{prop:twopass}. Can its compositions be evaluated directly
through the earlier component kernel, without lifting to the full dot basis?
Measure coefficient density and conversion costs, not just the nominal
factor-two dimension reduction. Prove that the selected marked point is
transported correctly through all local interfaces.
\end{question}

\begin{question}[A sharper support-dependent nilpotence bound]
The universal correction order is sharp, but most measured complexes use only
a small part of the arc algebra. Can one compute a smaller nilpotence bound for
the ideal generated by the actual entries of $\delta$, or for the operator
$h\delta$ on the survivor columns, without doing the transfer first? Intrinsic
weights give one inexpensive bound; radical-layer coordinates and forbidden
matching transitions may improve it.
\end{question}

\begin{question}[Transfer-compatible twist and separator macros]
Combine twist-normal-form complexes from the earlier package with this
survivor-first reduction. What contracts must a macro expose for its inclusion,
projection, and homotopy to remain compatible with subsequent gluing?
For separator-based recursion, bound both boundary width and minimal
multiplicity. A small number of geometric matching types by itself does not
control the number of copies of those types.
\end{question}

\begin{question}[Adaptive reduction with a provable safety envelope]
Can cheap predictors based on $(M,R,k)$, scalar rank density, and observed
composition cost select transfer or scalar cancellation reliably? A full-run
interleaving of two exact backends gives the usual constant-factor simulation
bound relative to the faster run when preemption overhead is bounded. Local
stage-by-stage switching is subtler: different minimal bases can alter the
cost of future stages even when their object profiles coincide. Establish the
appropriate claim rather than assuming local optimality implies global
optimality.
\end{question}

\begin{question}[Proof-carrying profile and contraction certificates]
Design serialized binary-rank and transfer certificates whose verification
avoids trusting the elimination implementation. The scalar contraction
identities already certify the predicted object counts, and
\cref{eq:certificate} certifies the transferred complex. Can certificates be
shared across many stages, with explicit bit-size bounds and without expanding
all homotopy matrices? A compressed certificate requires a correspondingly
proved verifier, not just a smaller file.
\end{question}

\begin{question}[Formalization of the characteristic-two core]
Formalize the finite matching category, the dot derivation, weight additivity,
and the sharp radical bound in Lean. A reusable second layer should prove
finite homological perturbation in an additive category with a nilpotent ideal.
The repository-specific interface can then instantiate these abstractions.
Checking $d^2=0$ is necessary but does not establish that a generated complex
represents the intended knot; local tangle equivalences remain part of the
formalization target.
\end{question}

\begin{question}[Which abstract sharpness witnesses occur topologically?]
The zigzag witnesses attain correction order $2k-2$ in the algebraic category.
Determine which can occur as subcomplexes or retracts of genuine knot-tangle
complexes. A topological restriction excluding the worst algebraic witnesses
could justify a better scanner-specific bound; a realizability construction
would rule out that particular improvement. Neither conclusion follows from
the finite algebra tests alone.
\end{question}

\begin{question}[A decision-specific interface invariant]
Can a smaller invariant than the full minimal complex decide whether an
arbitrary future closure has reduced Khovanov rank one, perhaps for a restricted
family of closure contexts? Universal composition summaries already face
strong dimensional constraints. A decision-only summary may be smaller, but
its closure domain and exactness must be stated, and positive and negative
answers must both be justified.
\end{question}

\begin{question}[Separating algebraic and hierarchy complexity]
The hierarchy program needs compressed cutting, controlled encodings, and
quantitative bounds on its geometric routines. Can the present exact rank and
nilpotence diagnostics serve as auxiliary verifiers or local cost estimates
without replacing those missing topological arguments? Any synthesis should
state a complete dependency graph and charge for every conversion between a
hierarchy description and an algebraic tangle description.
\end{question}

\section{Conclusion}
The practical improvement is a change in the order of work: determine the exact
minimal object profile with binary algebra first, then transfer the expensive
radical differential directly to those survivors. The accompanying structural
result is the sharp characteristic-two radical index $2k$, together with its
marked splitting and the sharp transfer order $2k-2$. These bounds are
independent of matrix multiplicity; that is precisely why they are useful for
large blocks of repeated matching objects.

The measurements support keeping both reduction strategies. They also support
integrating the survivor oracle before making more invasive changes. The
remaining global obstacle is no longer hidden inside an elimination heuristic:
it is the need to control, avoid, or succinctly represent the exact minimal
multiplicities and to find suitable presentations within a proved time bound.
The report supplies a parameterized theorem and a concrete model-cover target,
not an unrestricted quasi-polynomial recognizer.

\appendix
\section{A six-object example where the last correction is essential}
For $k=2$, let $a=(01)(23)$ and $b=(03)(12)$. Let
$f:a\to b$ and $g:b\to a$ be the undotted morphisms. Their composite is
$gf=x_1+x_2$ in $\End(a)$. Take degree-zero objects $(U,B_1,B_2)$ of types
$(a,b,a)$ and degree-one objects $(C_1,C_2,V)$ of types $(b,a,a)$, with
\[
 d=\begin{pmatrix}
 f&1&0\\
 0&g&1\\
 0&0&x_1
 \end{pmatrix}.
\]
The scalar residue consists only of the two identity pivots. Its survivors are
$U$ and $V$. The degree-zero and degree-one labels are different copies even
when they have the same matching type.

The direct term $p\delta i$ vanishes. The one-homotopy term
$p\delta h\delta i$ also vanishes because it ends at $C_2$. The last allowed
term is
\[
 p\delta h\delta h\delta i=x_1gf=x_1(x_1+x_2)=x_1x_2\ne0.
\]
Thus replacing the original complex by two isolated objects would be wrong,
even though it would have the correct minimal object profile. The exact
transferred differential is the fully dotted map between them.

\section{Source audit and reproducibility map}
\label{app:audit}
The following Git object identifiers specify the inspected upstream versions.
A hash identifies content; it is not a correctness certificate or an assertion
that an entire checkout was tested.
\begin{center}\footnotesize
\begin{tabular}{p{0.39\linewidth}p{0.55\linewidth}}
\toprule
Inspected source & Git object identifier\\\midrule
ProveIt \code{fast/} subtree & \code{b99a666269f660d932fc82ac8d1f29d3ae42231f}\\
\code{fastunknot/scan.py} & \code{4acff3cb5ffb4c4ad226a8bf6dcda3d8dc0bb561}\\
\code{fastunknot/algebra.py} & \code{d6d28604ac39c29145aa52f91074539a5aeac949}\\
\code{fastunknot/geometry.py} & \code{96acc160dbf2df7981b2db98b602311712a01531}\\
\code{fastunknot/planar.py} excerpt provenance & \code{1078526e7e7dbaf0b7267728105d870d94136769}\\
\code{openai/math} binary-algebra section & \code{29714bb7882fd989f987423540541957d8be6200}\\
\bottomrule
\end{tabular}
\end{center}

\noindent The archive is organized as follows.
\begin{center}\small
\begin{tabular}{p{0.37\linewidth}p{0.57\linewidth}}
\toprule
Path & Role\\\midrule
\code{article/article.tex}, \code{article.pdf} & Complete source and compiled report\\
\code{code/radical.py} & Exact residue, contraction, transfer, marked variant, pivot baseline\\
\code{code/braid_scan.py} & Crossing-attachment and closure validation harness\\
\code{vendor/reference_planar.py} & Attributed static numerical composition excerpt\\
\code{vendor/reference_cube.py} & Independent reduced crossing-cube oracle from prior package\\
\code{tests/} & Unit tests and deterministic complex fixtures\\
\code{experiments/} & Algebra, braid, marked-variant, and paired timing scripts\\
\code{results/} & Completed-run data and logs\\
\code{integration/} & Non-mutating reference-interface adapter and deployment notes\\
\code{PROVENANCE.json}, \code{CLAIMS.md} & Source identity and exact claim boundaries\\
\code{SHA256SUMS} & Hash manifest of delivered files\\
\bottomrule
\end{tabular}
\end{center}

The recommended reproduction commands are
\begin{lstlisting}
python -m unittest discover -s tests -v
python experiments/validate_algebra.py
python experiments/validate_braids.py
python experiments/validate_marked.py
python experiments/benchmark.py
python experiments/make_tables.py
cd article
pdflatex -halt-on-error article.tex
pdflatex -halt-on-error article.tex
\end{lstlisting}
The benchmark script records fresh timings, which naturally need not reproduce
identical milliseconds. The table generator keeps the numerical tables and
summary macros consistent with the recorded benchmark file. The SHA-256
manifest covers the delivered snapshot; a local rerun changes outputs and
therefore requires a new manifest to describe that new snapshot.

\begin{thebibliography}{99}
\bibitem{khovanov}
M. Khovanov, \emph{A functor-valued invariant of tangles}, Algebraic \& Geometric
Topology \textbf{2} (2002), 665--741.
\url{https://doi.org/10.2140/agt.2002.2.665}.
\bibitem{bn05}
D. Bar-Natan, \emph{Khovanov's homology for tangles and cobordisms}, Geometry \&
Topology \textbf{9} (2005), 1443--1499.
\url{https://doi.org/10.2140/gt.2005.9.1443}.
\bibitem{bn07}
D. Bar-Natan, \emph{Fast Khovanov homology computations}, Journal of Knot Theory
and Its Ramifications \textbf{16} (2007), 243--255.
\url{https://doi.org/10.1142/S0218216507005294}.
\bibitem{crainic}
M. Crainic, \emph{On the perturbation lemma, and deformations},
arXiv:math/0403266v1 (2004), especially Sections~2.2--2.4 and 3.2.
\url{https://arxiv.org/abs/math/0403266}.
\bibitem{km}
P. B. Kronheimer and T. S. Mrowka, \emph{Khovanov homology is an unknot-detector},
Publications Math\'ematiques de l'IH\'ES \textbf{113} (2011), 97--208.
\url{https://doi.org/10.1007/s10240-010-0030-y}.
\bibitem{ks26}
T. Kelom\"aki and D. Sch\"utz, \emph{On computational complexity of Khovanov
homology}, arXiv:2601.02119v1 (2026), especially Theorems~1.2 and 1.3.
\url{https://arxiv.org/abs/2601.02119}.
\bibitem{lackenby}
M. Lackenby, \emph{Incompressible surfaces, hierarchies and unknot recognition},
arXiv:2607.23350v1 (2026), Section~9.
\url{https://arxiv.org/abs/2607.23350}.
\bibitem{proveit}
V. Reshetnikov, \emph{ProveIt: Topology/UnknotRecognition}, repository sources
inspected 7 October 2026. Specific content identifiers appear in
\cref{app:audit} and the artifact provenance file.
\url{https://github.com/VladimirReshetnikov/ProveIt/tree/main/Topology/UnknotRecognition}.
\bibitem{mathrepo}
OpenAI, \emph{Mathematical manuscripts and supporting proof artifacts},
\code{openai/math}. The inspected methodological source is
\emph{Goldfeld's analytic density conjecture and the 2-converse for elliptic
curves}, binary-algebra section, subsection ``Finite cochain models,''
September 2026. No arithmetic theorem from that manuscript is used here.
\url{https://github.com/openai/math}.
\bibitem{twist}
ProveIt research contribution, \emph{Twist-Compressed Khovanov Recognition: Explicit complexity bounds,
exact cost certificates, and a tested route toward quasi-polynomial time}, 7 October 2026. Prior project artifact
\code{unknot_twist_research_bundle.zip}; exact source artifact and hashes are
recorded in \code{PROVENANCE.json}. Used for context and the independent
crossing-cube validation oracle.
\bibitem{component}
ProveIt research contribution, \emph{Component Compression and Restart Entropy
in Exact Unknot Recognition}, 7 October 2026. Prior project artifact
\code{unknot_component_entropy_20261007.zip}; exact source artifact and hashes
are recorded in \code{PROVENANCE.json}. Used for context and the attributed
static composition excerpt.
\end{thebibliography}
\end{document}
